\section{Problem class, decompositions and conditioning}\label{sec:setting}

\subsection{Model}
We study
\begin{equation}\label{eq:model}
\min_{x\in X}\ F(x)=\sum_{a\in\mathcal A}f_a(x_{S_a}),\qquad
X=\prod_{i=1}^nX_i ,
\end{equation}
where $X_i=[\ell_i,u_i]$ for $i\in I_C$ (continuous coordinates) and
$X_i=[\ell_i,u_i]\cap\Z$ for $i\in I_Z$ (integer coordinates). Integer
endpoints are rounded inward, an empty domain makes the problem infeasible,
and a coordinate with $\ell_i=u_i$ is substituted out; so we assume
$\ell_i<u_i$ for all $i$, with $\ell_i,u_i\in\Z$ for $i\in I_Z$. We write
$s_i=u_i-\ell_i$, $s=\max_is_i$, and $\bar X=\prod_i[\ell_i,u_i]$ for the
\emph{continuous hull} of $X$. Each factor $f_a$ depends only on the
coordinates in its scope $S_a\subseteq[n]=\{1,\dots,n\}$ and is defined and
continuous on the projection of $\bar X$, including fractional values of
integer coordinates. We write $\OPT=\min_XF$ and $\mathcal S$ for the set of
global minimizers.

\paragraph{Tree decompositions.}
A tree decomposition of the scopes is a tree $T$ with node set $\mathcal T$,
$N=\abs{\mathcal T}$, and bags $B_t\subseteq[n]$ such that every scope $S_a$
is contained in some bag and, for each coordinate $i$, the nodes whose bags
contain $i$ form a nonempty connected subtree. Each factor is assigned to
one bag containing its scope and each coordinate to one \emph{home bag}
containing it. We write $p=\max_t\abs{B_t}$ for the maximum bag size, so the
width of the decomposition is $p-1$. The decomposition is part of the input;
its validity is checked directly. A decomposition of width at most
$2\,\mathrm{tw}+1$, where $\mathrm{tw}$ is the treewidth of the interaction
graph, can be computed in time $2^{O(\mathrm{tw})}n$ \cite{Korhonen2021}, and
good heuristic decompositions are standard \cite{BodlaenderKoster2010}. No bound
on the number of bags containing a coordinate is assumed.

\paragraph{Encoding.}
For complexity statements the data are rational: box endpoints, and factors
given as explicit lists of monomials with rational coefficients. The input
length $I$ is the binary length of these data, including the decomposition
and the assignment of factors to bags. Running times are counted in bit
operations, and a target accuracy is written $2^{-q}$ with $q\in\Z_{\ge0}$.

\subsection{Coordinate curvature}
The certificate of Section~\ref{sec:grids} needs only one-dimensional upper
curvature information.

\begin{definition}[Coordinate curvature]\label{def:curvature}
Numbers $L_1,\dots,L_n\ge0$ are \emph{upper coordinate curvatures} of $F$ if,
for every $i$ and every $x\in\bar X$, the function
$t\mapsto F(x+te_i)-\frac{L_i}{2}t^2$ is concave on
$\{t:x+te_i\in\bar X\}$. We put $P=\{i:L_i>0\}$, $n_P=\abs P$, $L=\max_iL_i$,
and
\[
\norm d_{\mathbf L}^2=\sum_{i\in P}L_id_i^2\qquad(d\in\R^n).
\]
\end{definition}

The condition concerns the sum $F$, not the individual factors, and it must
hold on the continuous hull, including fractional points between integer
labels. If $F$ is twice differentiable it says $\partial_{ii}F\le L_i$ on
$\bar X$. For a quadratic $F(x)=\frac12x^\T Hx+b^\T x+c$ the smallest valid
choice is $L_i=\max\{H_{ii},0\}$; the off-diagonal entries of $H$ are
unrestricted, so $F$ can be far from convex. For explicit polynomials, valid
rational bounds follow from interval evaluation of $\partial_{ii}F$ over $\bar X$
(Lemma~\ref{lem:intcurv}). A coordinate with $L_i=0$ is concave along its own
direction, so replacing $x_i$ by one of $\ell_i,u_i$ never increases $F$; this
classical endpoint property \cite{Rosenberg1972,DelPiaKhajavirad2026} is the
case of zero correction below.

\subsection{Quadratic growth and conditioning}
The running-time analysis uses a growth condition at a unique minimizer.

\begin{definition}[Growth and condition numbers]\label{def:growth}
Problem~\eqref{eq:model} has \emph{quadratic growth} with constant $g>0$ if
there is $x^*\in X$ with
\begin{equation}\label{eq:growth}
F(x)-F(x^*)\ge g\norm{x-x^*}^2\qquad\text{for all }x\in X ;
\end{equation}
then $x^*$ is the unique minimizer, and we put $\kappa=\max\{1,L/g\}$. It has
\emph{weighted growth} with constant $\gamma>0$ if there is a minimizer $x^*$ with
\begin{equation}\label{eq:wgrowth}
F(x)-\OPT\ge\gamma\norm{x-x^*}_{\mathbf L}^2\qquad\text{for all }x\in X .
\end{equation}
Let $\gamma^*$ be the largest such constant and $\bar\kappa=\max\{1,1/\gamma^*\}$
($\bar\kappa=\infty$ if there is none).
\end{definition}

Since $\norm d_{\mathbf L}^2\le L\norm d^2$, quadratic growth with constant
$g$ implies weighted growth with constant $g/L$, so $\bar\kappa\le\kappa$. The
weighted condition number has two advantages. It does not change when a
continuous coordinate is rescaled, $x_i=\alpha y_i$, because $L_i$ is multiplied
by $\alpha^2$ and $L_i(x_i-x_i^*)^2$ is unchanged; for a separable quadratic
$\sum_i\frac{L_i}{2}(x_i-c_i)^2$ with interior minimizer, $\bar\kappa=2$ while
$\kappa=2\max_iL_i/\min_iL_i$. And coordinates with $L_i=0$ need no growth at
all: weighted growth makes $x^*_P$ unique but allows several minimizers that
differ outside $P$.

Both conditions are global and imposed only on the feasible set $X$. They do
not require convexity: $F$ may have many local minima, provided each lies above
$\OPT$ by an amount proportional to its squared distance from $x^*$
(Example~\ref{ex:family}). When integer coordinates are present, nearly tied
integer assignments make the growth constant small even if $F$ is strongly
convex in the continuous coordinates. For quadratics, uniqueness of the
minimizer already implies~\eqref{eq:growth} for some $g>0$
(Lemma~\ref{lem:unique-growth}), but $g$ can be exponentially small in $I$.
Our algorithms use the curvature bounds explicitly; they never need $g$,
$\gamma$ or $x^*$, and the validity of every bound they return depends only on
Definition~\ref{def:curvature}.

For box QPs, the next lemma certifies growth from the gradient and the
Hessian at a point, and shows what minimality and growth imply for the
Hessian. For a point $x^*\in X$ let
$J_0=\{i\in I_C:\ell_i<x^*_i<u_i\}$ be its continuous coordinates strictly
inside their bounds and $J_\partial=\{i:x^*_i\in\{\ell_i,u_i\}\}$ its
coordinates at a bound.

\begin{lemma}[Growth certificates for box QPs]\label{lem:growthcert}
Let $F=\frac12x^\T Hx+b^\T x+c$ be a quadratic on a mixed box $X$, $x^*\in X$
and $\zeta=\nabla F(x^*)$.
\begin{enumerate}
\item[(a)] Suppose that $\zeta_i=0$ for $i\in J_0$, and that for every
$i\in I_C\setminus J_0$ the gradient points inward: $\zeta_i\ge0$ if
$x^*_i=\ell_i$ and $\zeta_i\le0$ if $x^*_i=u_i$. Let $M=\diag(\mu)$ with
$\mu_i=0$ for $i\in J_0$, $\mu_i=\abs{\zeta_i}/s_i$ for $i\in J_\partial$
with inward $\zeta_i$, and $\mu_i=-\abs{\zeta_i}$ for the remaining
coordinates, which are integer coordinates. Then
\begin{equation}\label{eq:minorant}
F(x)-F(x^*)\ge\tfrac12(x-x^*)^{\T}(H+2M)(x-x^*)\qquad(x\in X).
\end{equation}
Consequently, if $H+2M\succeq2g_0I$ for some $g_0>0$, then $x^*$ is the unique
minimizer and $F$ has point growth with constant $g_0$.
\item[(b)] Conversely, let $x^*$ be a minimizer. Then $\zeta_{J_0}=0$ and
$H_{J_0J_0}\succeq0$. If $F$ has point growth with constant $g$, then
$H_{J_0J_0}\succeq2gI$; if it has weighted growth with constant $\gamma$ at
$x^*$, then $H_{J_0J_0}\succeq2\gamma\diag(L_i)_{i\in J_0}$. In particular,
under point growth $g\le\frac12v^{\T}Hv/\norm v^2$ for every nonzero $v$
supported on $J_0$, and if $J_0\ne\emptyset$, every valid curvature bound
satisfies $L\ge2g$, so $\kappa\ge2$. If $J_0=[n]$, then $F$ is convex, and
strongly convex under point growth.
\end{enumerate}
\end{lemma}

\begin{proof}
(a) Let $x\in X$ and $d=x-x^*$. Since $F$ is quadratic,
$F(x)-F(x^*)=\zeta^{\T}d+\frac12d^{\T}Hd$ exactly, and it suffices to show
$\zeta_id_i\ge\mu_id_i^2$ for every $i$. For $i\in J_0$ both sides vanish. If
$x_i^*=\ell_i$ and $\zeta_i\ge0$, then $0\le d_i\le s_i$, so
$\zeta_id_i=\abs{\zeta_i}d_i\ge\abs{\zeta_i}d_i^2/s_i=\mu_id_i^2$; the case
$x_i^*=u_i$, $\zeta_i\le0$ is symmetric. For the remaining $i\in I_Z$, $d_i\in\Z$
and $\abs{d_i}\le d_i^2$, which gives
$\zeta_id_i\ge-\abs{\zeta_i}\abs{d_i}\ge-\abs{\zeta_i}d_i^2=\mu_id_i^2$. If
$H+2M\succeq2g_0I$, the right side of~\eqref{eq:minorant} is at least
$g_0\norm d^2$, which is positive for $d\ne0$.

(b) For $i\in J_0$ both directions $\pm e_i$ are feasible from $x^*$, so
minimality gives $\zeta_i=0$. Let $v\ne0$ be supported on $J_0$. For all
sufficiently small $\abs t$, $x^*+tv\in X$ and
$F(x^*+tv)-\OPT=\frac{t^2}2v^{\T}Hv$. This is nonnegative by minimality;
point growth bounds it from below by $gt^2\norm v^2$, and weighted growth by
$\gamma t^2\norm v_{\mathbf L}^2=\gamma t^2\sum_{i\in J_0}L_iv_i^2$; this gives
the three matrix inequalities. Taking $v=e_i$ with $i\in J_0$ gives
$H_{ii}\ge2g$ under point growth, and every valid curvature bound satisfies
$L\ge L_i\ge H_{ii}$, because $t\mapsto F(x+te_i)-L_it^2/2$ has second
derivative $H_{ii}-L_i$. If $J_0=[n]$, then $H=H_{J_0J_0}$.
\end{proof}

\begin{remark}[What growth restricts]\label{rem:nonconvex}
For a box QP, the block $H_{J_0J_0}$ at a minimizer is positive semidefinite
(Lemma~\ref{lem:growthcert}(b)), so at every minimizer, with or without
growth, every direction of negative curvature involves active bounds or
integer coordinates. Point growth with some $g>0$ is equivalent to
uniqueness of the minimizer (Lemma~\ref{lem:unique-growth}), so what
restricts the instances is the size of the constants. Point growth gives
$H_{J_0J_0}\succeq2gI$, hence $\kappa\ge2L/\lambda_{\min}(H_{J_0J_0})$ if
$J_0\ne\emptyset$, and weighted growth gives
$H_{J_0J_0}\succeq2\gamma\diag(L_i)_{i\in J_0}$. For every $x'\in X$ with
$F(x')>\OPT$, \eqref{eq:growth} gives
$\kappa\ge L\norm{x'-x^*}^2/(F(x')-\OPT)$ and \eqref{eq:wgrowth} gives
$\bar\kappa\ge\norm{x'-x^*}_{\mathbf L}^2/(F(x')-\OPT)$. Since
$\norm d_{\mathbf L}^2\le L\norm d^2$, nearly optimal points far from $x^*$
in the norm $\norm\cdot_{\mathbf L}$ make both condition numbers large.
Within these conditions an instance can have an indefinite Hessian,
$\kappa\le\frac{21}4$ and exponentially many strict local minima
(Example~\ref{ex:family}).
Part~(a) certifies growth when the first-order terms $\zeta_id_i$ at active
bounds supply the growth that the Hessian lacks. At integer coordinates it
uses integrality only to bound $\zeta_id_i$ below by $-\abs{\zeta_i}d_i^2$.
\end{remark}

